% Valentino Vranic
% Metody inzinierskej prace 2012/13

\documentclass{beamer}

%\usetheme{Warsaw}
\usetheme{Antibes}
%\usetheme{JuanLesPins}
%\usetheme{Goettingen}

%\usecolortheme{seahorse}
%\usecolortheme{dolphin}
%\usecolortheme{rose}
% http://deic.uab.es/~iblanes/beamer_gallery/index_by_color.html
%\usecolortheme{beaver}
\usecolortheme{orchid}

%\useoutertheme[]{sidebar}

\setbeamercovered{transparent}

\usepackage[slovak]{babel}
\usepackage[T1]{fontenc}
\usepackage[utf8]{inputenc}
\usepackage{url}

\usepackage{array}

\usepackage{listings}

\lstset{language=Python,basicstyle=\fontsize{8}{9.6}\selectfont,showstringspaces=false,columns=fullflexible,identifierstyle=\ttfamily,keywordstyle=\bfseries,showstringspaces=false,columns=fullflexible}
%\lstset{language=C,basicstyle=\fontsize{10.5}{12.6}\selectfont,identifierstyle=\ttfamily,keywordstyle=\bfseries,showstringspaces=false,columns=fixed}

\def\BibTeX{\textsc{Bib}\kern-.08em\TeX} 

\newcommand{\footcite}[1]{\footnote{\tiny #1}}
\newcommand{\umlet}{.5}
\newcommand{\emp}[1]{\textit{\alert{#1}}}
\newcommand{\kw}[1]{\mbox{\textbf{#1}}}
\newcommand{\id}[1]{\texttt{#1}}
\newcommand{\stl}{\guillemotleft}
\newcommand{\str}{\guillemotright}

\newcommand{\lsti}{\lstinline[basicstyle=\fontsize{10.5}{12.1}\selectfont]}

\newcommand{\ssection}[1]{
	\section{#1}
	\begin{frame}[fragile=singleslide]\frametitle{}
	\Huge #1
	\end{frame}
}

\newcommand{\ssectionn}[1]{
	\section*{#1}
	\begin{frame}[fragile=singleslide]\frametitle{}
	\Huge #1
	\end{frame}
}

\newenvironment{program}{\begin{beamercolorbox}[rounded=true,shadow=true]{block body}\vspace{-4mm}}{\vspace{-2mm}\end{beamercolorbox}}

\setbeamercolor{fvystup}{fg=white,bg=black}
\newenvironment{vystup}{\begin{beamercolorbox}[rounded=true,shadow=true]{fvystup}}{\end{beamercolorbox}}

\newenvironment{poznamka}{\begin{beamercolorbox}[rounded=true,shadow=false]{block body}}{\end{beamercolorbox}}

\setbeamertemplate{footline}[page number]
{
%\insertpagenumber
%\begin{beamercolorbox}{section in head/foot}
%\vskip2pt\insertnavigation{\paperwidth}\vskip2pt
%\end{beamercolorbox}%
}



\author{Štefan Kučerák\\{\tiny\url{xkucerak@stuba.sk}}}
%\url{www.fiit.stuba.sk/~vranic}, \url{vranic@fiit.stuba.sk}}
%{\tiny \url{www.fiit.stuba.sk/~vranic}, \url{vranic@fiit.stuba.sk}}
\institute{
	Ústav informatiky, informačných systémov a softvérového inžinierstva\\
	Fakulta informatiky a informačných technológií\\
	Slovenská technická univerzita v Bratislave}

\subtitle{\vspace{3mm} Metódy inžinierskej práce 2023/2024}

\title{Extrakcia informácii z webu pre analýzu dát pomocou Pythonu}

\date{\footnotesize 27. november 2023}




\begin{document}

\begin{frame}[fragile=singleslide]
\titlepage
\end{frame}


\begin{frame}[fragile=singleslide]\frametitle{Introduction}
Web scraping:
\begin{itemize}
\item automated data collection from websites
\item obtaining a large amount of data in a short time 
\end{itemize}

\end{frame}

\begin{frame}[fragile=singleslide]\frametitle{Overview}
\tableofcontents
\end{frame}

\section{Principle of data extraction}
\begin{frame}[fragile=singleslide]\frametitle{The principle of data extraction}
\includegraphics[scale=0.75]{diag.png}
% pridajte vlastný obrázok a zrušte znák % pred príkazom \includegraphics vo formáte PDF prípadne PNG alebo JPG
% scale určuje veľkosť obrázku

{\tiny Principle of data extraction.}
\end{frame}

\begin{frame}[fragile=singleslide]\frametitle{Libraries}
\begin{center}
\begin{table}[]
\begin{tabular}{ l | p{2.5in} }
Library& Description\\\hline
Requests\footcite{\url{https://pypi.org/project/requests/}} & HTTP requests using POST and GET\\
Beautiful Soup\footcite{\url{https://pypi.org/project/beautifulsoup4/}} & the library contains tools for viewing HTML or XML code\\
Scrapy\footcite{\url{https://pypi.org/project/Scrapy/}} & is a fast framework for high-level web browsing\\
Selenium\footcite{\url{https://pypi.org/project/selenium/}} & is a library that is a web driver\end{tabular}
\end{table}
\end{center}
\end{frame}



\section{Exmaples of use in real life}
% príkaz \ssection by vytvoril zvláštný slajd s názvom časti - v krátkych prezentáciách to prekáža, lebo oberá o čas

\begin{frame}[fragile=singleslide]\frametitle{Examples}
The following projects used Python and one of the libraries listed: 
\begin{itemize}
\item Phishing Web Page Detection using web Scraping
\item Food Genie, Recipe Search Algorithm using Web Scraping
\item NEWSONE- AN AGGREGATION SYSTEM FOR NEWS USING WEB SCRAPING METHOD
\end{itemize}
\end{frame}


\section{Example code}
\begin{frame}[fragile=singleslide]\frametitle{Simple example}
\begin{poznamka}
Listing the page title using the Beautiful Soup library
\end{poznamka}

\begin{itemize}
\item CODE

\begin{program}
\begin{lstlisting}
import requests
from bs4 import BeautifulSoup

page = requests.get("https://www.google.com/")
soup = BeautifulSoup(page.content, "html.parser")

print(soup.title.string)
\end{lstlisting}
\end{program}

\item Output
\begin{vystup}
\begin{lstlisting}
Google
\end{lstlisting}
\end{vystup}

\end{itemize}
\end{frame}

\begin{frame}[fragile=singleslide]\frametitle{This solution can be used on multiple sites}

\includegraphics[scale=0.45]{tabulka.png}

\end{frame}



\section*{Conclusion}
% hviezdička zabezpečí, aby sa táto časť neocitla v prehľade prezentácie - každá prezentácia má zhodnotenie a prehľad by sa tým zbytočne zahlcoval

\begin{frame}[fragile=singleslide]\frametitle{Conclusion}
\begin{itemize}
\item the Python programming language is widely used in this area
\item we can use one of the listed libraries to extract data from the web 
\item we can then use the data for further processing 
\end{itemize}
\end{frame}


\end{document}




Text \end{document} za príkazom \end{document} LaTeX ignoruje, takže tu môžete odkladať veci (aj celé slajdy), ktoré nechcete vymazať, lebo ich ešte možno budete potrebovať, avšak ich v danom momente nechcete mať v slajdoch.
